\chapter*{TO THE READER}
\phantomsection
\addcontentsline{toc}{chapter}{To the Reader}

\begin{enumerate}
\def\labelenumi{\arabic{enumi}.}
\item
  This series of volumes, a list of which is given on pages vii and viii, takes up mathematics at the beginning, and gives complete proofs. In principle, it requires no particular knowledge of mathematics on the reader's part, but only a certain familiarity with mathematical reasoning and a certain capacity for abstract thought. Nevertheless, it is directed especially to those who have a good knowledge of at least the content of the first year or two of a university mathematics course.
\item
  The method of exposition we have chosen is axiomatic and abstract, and normally proceeds from the general to the particular. This choice has been dictated by the main purpose of the treatise, which is to provide a solid foundation for the whole body of modern mathematics. For this it is indispensable to become familiar with a rather large number of very general ideas and principles. Moreover, the demands of proof impose a rigorously fixed order on the subject matter. It follows that the utility of certain considerations will not be immediately apparent to the reader unless he has already a fairly extended knowledge of mathematics; otherwise he must have the patience to suspend judgment until the occasion arises.
\item
  In order to mitigate this disadvantage we have frequently inserted examples in the text which refer to facts the reader may already know but which have not yet been discussed in the series. Such examples are always placed between two asterisks: * \ldots{} *. Most readers will undoubtedly find that these examples will help them to understand the text, and will prefer not to leave them out, even at a first reading. Their omission would of course have no disadvantage, from a purely logical point of view.
\item
  This series is divided into volumes (here called ``Books''). The first six Books are numbered and, in general, every statement in the text assumes as known only those results which have already been discussed in the preceding volumes. This rule holds good within each Book, but for convenience of exposition these Books are no longer arranged in a consecutive order. At the beginning of each of these Books (or of these chapters), the reader will find a precise indication of its logical relationship to the other Books and he will thus be able to satisfy himself of the absence of any vicious circle.
\item
  The logical framework of each chapter consists of the definitions, the axioms, and the theorems of the chapter. These are the parts that have mainly to be borne in mind for subsequent use. Less important results and those which can easily be deduced from the theorems are labelled as ``propositions'', ``lemmas'', ``corollaries'', ``remarks'', etc. Those which may be omitted at a first reading are printed in small type. A commentary on a particularly important theorem appears occasionally under the name of ``scholium''.
\end{enumerate}

To avoid tedious repetitions it is sometimes convenient to introduce notations or abbreviations which are in force only within a certain chapter or a certain section of a chapter (for example, in a chapter which is concerned only with commutative rings, the word ``ring'' would always signify ``commutative ring''). Such conventions are always explicitly mentioned, generally at the beginning of the chapter in which they occur.

\begin{enumerate}
\def\labelenumi{\arabic{enumi}.}
\setcounter{enumi}{5}
\tightlist
\item
  Some passages in the text are designed to forewarn the reader against serious errors. These passages are signposted in the margin with the sign
\end{enumerate}

Z (``dangerous bend'').

\begin{enumerate}
\def\labelenumi{\arabic{enumi}.}
\setcounter{enumi}{6}
\item
  The Exercises are designed both to enable the reader to satisfy himself that he has digested the text and to bring to his notice results which have no place in the text but which are nonetheless of interest. The most difficult exercises bear the sign ♦.
\item
  In general, we have adhered to the commonly accepted terminology, except where there appeared to be good reasons for deviating from it.
\item
  We have made a particular effort always to use rigorously correct language, without sacrificing simplicity. As far as possible we have drawn attention in the text to abuses of language, without which any mathematical text runs the risk of pedantry, not to say unreadability.
\item
  Since in principle the text consists of the dogmatic exposition of a theory, it contains in general no references to the literature. Bibliographical references are gathered together in Historical Notes, usually at the end of each chapter. These notes also contain indications, where appropriate, of the unsolved problems of the theory.
\end{enumerate}

The bibliography which follows each historical note contains in general only those books and original memoirs which have been of the greatest importance in the evolution of the theory under discussion. It makes no sort of pretence to completeness; in particular, references which serve only to determine questions of priority are almost always omitted.

As to the exercises, we have not thought it worthwhile in general to indicate their origins, since they have been taken from many different sources (original papers, textbooks, collections of exercises).

\begin{enumerate}
\def\labelenumi{\arabic{enumi}.}
\setcounter{enumi}{10}
\tightlist
\item
  References to a part of this series are given as follows :
\end{enumerate}

\begin{enumerate}
\def\labelenumi{\alph{enumi})}
\item
  If reference is made to theorems, axioms, or definitions presented in the same section, they are quoted by their number.
\item
  If they occur in another section of the same chapter, this section is also quoted in the reference.
\item
  If they occur in another chapter in the same Book, the chapter and section are quoted.
\item
  If they occur in another Book, this Book is first quoted by its title.
\end{enumerate}

The Summaries of Results are quoted by the letter R: thus Set Theory, R signifies ``Summary of Results of the Theory of Sets''.

\chapter*{CONTENTS OF THE ELEMENTS OF MATHEMATICS SERIES}
\phantomsection
\addcontentsline{toc}{chapter}{Contents of the Elements of Mathematics Series}

\listheading{I. THEORY OF SETS}

\begin{enumerate}
\def\labelenumi{\arabic{enumi}.}
\tightlist
\item
  Description of formal mathematics. 2. Theory of sets. 3. Ordered sets; cardinals; natural numbers. 4. Structures.
\end{enumerate}

\listheading{II. ALGEBRA}

\begin{enumerate}
\def\labelenumi{\arabic{enumi}.}
\tightlist
\item
  Algebraic structures. 2. Linear algebra. 3. Tensor algebras, exterior algebras, symmetric algebras. 4. Polynomials and rational fractions. 5. Fields. 6. Ordered groups and fields. 7. Modules over principal ideal rings. 8. Semi-simple modules and rings. 9. Sesquilinear and quadratic forms.
\end{enumerate}

\listheading{III. GENERAL TOPOLOGY}

\begin{enumerate}
\def\labelenumi{\arabic{enumi}.}
\tightlist
\item
  Topological structures. 2. Uniform structures. 3. Topological groups. 4. Real numbers. 5. One-parameter groups. 6. Real number spaces, affine and projective spaces. 7. The additive groups \(\mathbb{R}^n\) . 8. Complex numbers. 9. Use of real numbers in general topology. 10. Function spaces.
\end{enumerate}

\listheading{IV. FUNCTIONS OF A REAL VARIABLE}

\begin{enumerate}
\def\labelenumi{\arabic{enumi}.}
\tightlist
\item
  Derivatives. 2. Primitives and integrals. 3. Elementary functions. 4. Differential equations. 5. Local study of functions. 6. Generalized Taylor expansions. The Euler-Maclaurin summation formula. 7. The gamma function. Dictionary.
\end{enumerate}

\listheading{V. TOPOLOGICAL VECTOR SPACES}

\begin{enumerate}
\def\labelenumi{\arabic{enumi}.}
\tightlist
\item
  Topological vector spaces over a valued field. 2. Convex sets and locally convex spaces. 3. Spaces of continuous linear mappings. 4. Duality in topological vector spaces. 5. Hilbert spaces: elementary theory. Dictionary.
\end{enumerate}

\listheading{VI. INTEGRATION}

\begin{enumerate}
\def\labelenumi{\arabic{enumi}.}
\tightlist
\item
  Convexity inequalities. 2. Riesz spaces. 3. Measures on locally compact spaces. 4. Extension of a measure. L \(^{p}\) spaces. 5. Integration of measures. 6. Vectorial integration. 7. Haar measure. 8. Convolution and representation.
\end{enumerate}

\listheading{LIE GROUPS AND LIE ALGEBRAS}

\listheading{1. Lie algebras.}

\listheading{COMMUTATIVE ALGEBRA}

\begin{enumerate}
\def\labelenumi{\arabic{enumi}.}
\tightlist
\item
  Flat modules. 2. Localization. 3. Graduations, filtrations and topologies. 4. Associated prime ideals and primary decomposition. 5. Integers. 6. Valuations. 7. Divisors.
\end{enumerate}

\listheading{SPECTRAL THEORIES}

\begin{enumerate}
\def\labelenumi{\arabic{enumi}.}
\tightlist
\item
  Normed algebras. 2. Locally compact groups.
\end{enumerate}

